\section{Day 3: Mary in the Heart of the Expectant Church}

\dayaxis{Luke 1:26--38; Acts 1:14; John 19:25--27}{Receptive faith, maternal intercession, and consent to divine initiative}{Pray for and perform one hidden act of care for the Church or a person entrusted to you.}

\subsection{Read: Luke 1:26--38 and Acts 1:14}

Acts names her with rare simplicity: “Mary the mother of Jesus.” That title keeps Pentecost joined to the Incarnation. The eternal Son truly took flesh from her; the glorified Christ remains her Son according to the humanity he assumed. Her divine motherhood is Christological before it is devotional.

At the Annunciation the Spirit's overshadowing is ordered to the conception of Christ. At Pentecost the Spirit's visible mission equips the apostolic Church for witness. These events should neither be collapsed nor isolated. Mary does not receive a second divine motherhood of the same order, and the apostles do not become incarnational mothers in her unique sense. Yet the pattern of grace and response endures: God initiates; his word is heard; a creature asks a real question without unbelief; grace enables consent; Christ is borne toward the world.

Mary's prayer is subordinate intercession within Christ's one mediation. \textit{Lumen gentium} 60--62 insists that no creature can be placed on the same level as the incarnate Word and Redeemer; Mary's maternal role depends entirely upon his merits and power. Precisely this protects the beauty of her presence. She need not rival the source of grace in order to be the Mother praying among Christ's members.

\subsection{Meditation: consent without possession}

“Be it done unto me according to thy word” is not passivity before evil or the surrender of reason. Mary has heard a divine message, received an answer, and consents to God's holy action. One must not use her fiat to pressure another person into abuse, unsafe silence, or an alleged private revelation. Christian surrender is to God as known through faith and right reason, never to sin.

Her presence also purifies the desire for prominence. Acts does not give her a speech, office-election, or self-explanation in the upper room. Hidden prayer can be fully fecund. The person who cannot lead a public ministry may intercede, suffer faithfully, teach a child, reconcile a family, support an apostolate, or guard the Church from gossip. The Spirit distributes roles; charity makes each one fruitful.

On a transferred-Ascension calendar, Day~3 may coincide with the liturgical solemnity. Let the day's Marian note remain transparent to Christ: the Mother looks toward her exalted Son and awaits his promise with his Church.

\novenaexamination{Do I call control “cooperation”? Do I despise hidden service because it is unseen? Is my Marian devotion clearly centered on Christ, the Trinity, Scripture, sacrament, conversion, and charity?}

\bilingualprayer{The Annunciation of the Blessed Virgin Mary}{Received 1962 Latin collect, with its abbreviated same-Lord conclusion mechanically completed from \textit{Rubricae generales} 115(a)--(b), p.~XIX; received historical English from Lefebvre's 1925 \textit{Daily Missal}, p.~1310, with its abbreviated same-Lord conclusion mechanically completed from that witness's full formula and response at p.~833. Historical source approval does not extend to this novena.}{%
Deus, qui de beátæ Maríæ Vírginis útero Verbum tuum, Angelo nuntiánte, carnem suscípere voluísti: præsta supplícibus tuis; ut, qui vere eam Genetrícem Dei crédimus, eius apud te intercessiónibus adiuvémur. Per eúndem Dóminum nostrum Iesum Christum, Fílium tuum: Qui tecum vivit et regnat in unitáte Spíritus Sancti, Deus, per ómnia s\'{\ae}cula sæculórum. Amen.}{%
O God Who didst please that Thy Word should take flesh, at the message of an Angel, in the womb of the blessed Virgin Mary, grant to Thy suppliants, that we who believe her to be truly the Mother of God, may be helped by her intercession with Thee. Through the same Lord Jesus Christ Thy Son: Who with Thee and the Holy Ghost liveth and reigneth one God for ever and ever. Amen.}


\prayerinstruction{Continue with \textit{Veni Creator Spiritus}, the Lord's Prayer, and the Lesser Doxology.}
